% Autor: Simon May
% Datum: 2014-08-12
\section[Die junge DPG stellt sich vor]{Mehr als nur Uni -- Die junge~DPG stellt sich vor}
\begin{minipage}{0.1\textwidth}
\includegraphics[width=\textwidth]{res/jdpg_logo.pdf}
\end{minipage}\hfill
\begin{minipage}{0.75\textwidth}
\centering
\textbf{Liebe Physik-Erstis der Uni Münster,\\
auch die junge DPG möchte Euch herzlich im Kreis der Physik-Studenten begrüßen.}
\end{minipage}\hfill
\begin{minipage}{0.1\textwidth}
\includegraphics[width=\textwidth]{res/jdpg_logo.pdf}
\end{minipage}

\begin{multicols}{2}
In Deutschland studieren rund \num{30000} Studenten Physik. Du bist jetzt einer von ihnen. An "Jungphysiker" wie dich wendet sich die junge~Deutsche~Physikalische~Gesellschaft~(jDPG) und macht ihnen Angebote rund um das Physikstudium.

Das Konzept der jDPG basiert auf fünf Säulen: Diese sind neben der Berufsvorbereitung und dem wissenschaftlichen Programm vor allem der Dialog zwischen den Physiker-Generationen und die Bildung eines Netzwerks für Studierende der Physik. Außerdem bildet die jDPG zusammen mit den Physikfachschaften eine bundesweite Interessenvertretung von Physikstudenten -- die fünfte Säule des Programms.

Der Kontakt zur jDPG ist leicht herzustellen und jeder ist willkommen. Am schnellsten findest du die Informationen unter

\begin{center}
\textbf{\url{http://www.jdpg.de}}.
\end{center}

Die Angebote der Regionalgruppe~Münster findest du auch direkt unter

\begin{center}
\textbf{\url{http://muenster.jdpg.de}}.
\end{center}

Mehr als \num{1200} Mitglieder zählt die jDPG derzeit. Für diese Mitglieder und meistens auch für alle anderen Physik-Studierenden gibt es unsere Angebote. Viele Regionalgruppen -- das sind Studentengruppen an einer bestimmten Uni -- veranstalten zum Beispiel das sogenannte "Meet your Prof". So auch die Regionalgruppe in Münster. Bei dieser Veranstaltungsreihe können Studenten ihre Professoren näher kennen lernen. In gemütlicher Atmosphäre -- mit Orangensaft und Salami-Brötchen -- sitzen die Studenten mit ihren Profs am Tisch. Diese erzählen von ihrer aktuellen Forschung, aber auch Fragen zum Privatleben sind möglich und sogar erwünscht. Das ist ein Teil des Dialogs, den die jDPG möchte: Studenten treffen Professoren, Professoren treffen Studenten.

Der Anstoß zur Gründung der jDPG kam im Jahr~2005 aus der Deutschen~Physikalischen~Gesellschaft selbst. Diese hatte zwar 30~Prozent studentische Mitglieder, aber die Angebote für diese junge Zielgruppe fehlten. Als die DPG auf die Studenten der TU~Dresden zuging, haben sich spontan fünf Jungphysiker gefunden, die bereit waren, diese Lücke zu schließen. "Wir haben das damals für ein kleines Versuchsprojekt gehalten", erzählt jDPG-Mitglied René Pfitzner. Mit der rasanten Entwicklung der vergangenen sechs Jahre habe damals niemand gerechnet.

Die erste Veranstaltung der jDPG war eine Sommerexkursion nach Berlin. Das Max-Born-Institut, das Albert-Einstein-Institut und auch die Firma Rolls~Royce standen auf dem Programm.

\begin{quote}
\textit{"Das hat uns einen richtigen Schub gegeben"},
\end{quote}
erinnert sich René Pfitzner. Aus ganz Deutschland hätten sich Freiwillige gemeldet, die die Idee der jDPG unterstützen wollten. So haben sich langsam die ersten Regionalgruppen gebildet und Ende~2006 wurde die jDPG offiziell gegründet.

Doch die steigenden Mitgliederzahlen zeigten auch, dass eine lose Arbeitsgruppe, wie sie am Anfang geplant war, nicht mehr ausreichen würde. Eine neue Struktur musste her und diese wurde auf der ersten Mitgliederversammlung im Jahr~2007 beschlossen und seitdem immer wieder verändert.

\begin{quote}
\textit{"Wir haben die jDPG nie auf dem Reißbrett geplant"},
\end{quote}
sagt der langjährige Bundessprecher Heinrich.
\begin{quote}
\textit{"Die Struktur ist immer aus den neuen Anforderungen gewachsen."}
\end{quote}

Heute ist die jDPG eine bundesweit aktive Organisation mit festem Programm. Vor Ort machen zurzeit 31~Regionalgruppen den Mitgliedern lokale Angebote.

Die Regionalgruppe Münster organisiert regelmäßig Exkursionen oder zeigt mit dem "EinBlick in den Beruf" an einem Abend Innenansichten eines bestimmten Arbeitsgebietes für Physiker.

Deutschlandweit ist die Sommerexkursion das Highlight im jDPG-Jahr. Die Idee ist noch dieselbe wie im Gründungsjahr~2005.
\begin{quote}
\textit{"Wir möchten Physikstudenten ihre unzähligen beruflichen Möglichkeiten aufzeigen und einen Einblick in aktuelle Forschung und Entwicklung geben"},
\end{quote}
erklärt Pfitzner.
\begin{quote}
\textit{"Und zwar nicht im Hörsaal, sondern direkt am Schauplatz des Geschehens."}
\end{quote}
Nach Bonn, Berlin, München, Heidelberg, Dresden und Hannover in den vergangenen Jahren ging die Reise 2013 nach Hamburg.

Ein ähnliches Ziel wie die "EinBlicke in den Beruf" verfolgen die Berufsvorbereitungsseminare~(BVS). Physiker aus verschiedenen Bereichen der Wirtschaft berichten über ihr Arbeitsgebiet und geben Hinweise für den Berufseinstieg. Auch hier legt die jDPG Wert auf eine lockere Atmosphäre und kleine Gruppen, sodass die Teilnehmer genug Möglichkeiten zum Fragen haben -- ob in der Gesprächsrunde am Vormittag oder bei einem Glas Wein am Abend. Das nächste BVS findet Ende November in Münster statt.

Wichtig für eine erfolgreiche Arbeit der jDPG ist laut Pfitzner auch eine effiziente Kommunikation mit den Mitgliedern. Dafür informiert die jDPG regelmäßig in einem eigenen Newsletter auf der Internetseite \url{http://www.jdpg.de} und im Physik~Journal, der Mitgliederzeitschrift der Deutschen~Physikalischen~Gesellschaft.

Zirka 60~aktive Mitglieder organisieren die Projekte der jDPG, aber ein Vorurteil will Anna Bakenecker (Bundessprecherin) schnell entkräften:
\begin{quote}
\textit{"Nicht jedes Mitglied muss aktiv mitarbeiten."}
\end{quote}
Viele Studenten würden dies falsch verstehen, da die jDPG offiziell eine Arbeitsgruppe der Deutschen Physikalischen Gesellschaft sei. "Arbeitsgruppe" habe aber mit Mitarbeit nichts zu tun. Anna Bakenecker sieht in der Werbung neuer aktiver Mitglieder einer der wichtigsten Aufgaben der kommenden Jahre. Das Tolle an der jDPG sei, dass man in kurzer Zeit viel erreichen könne, sagt die Physikstudentin.
\begin{quote}
\textit{"In vielen Organisationen versickern neue Ideen. Bei uns werden sie umgesetzt."}
\end{quote}

\includegraphicscompressed[width=\columnwidth]{res/jdpg_foto.png}

\fibelsig{Philipp Seibt und Markus Tegeder}
\end{multicols}

\begin{center}
\includegraphics[width=0.8\textwidth]{res/comics/xkcd_experiment.png}
\end{center}
